
\chapter{\textbf{INTRODUÇÃO}}
    Na Astrofísica contemporânea, a determinação precisa da massa estelar emerge como um desafio crucial. Esta compreensão detalhada é fundamental para a investigação da formação, evolução, dinâmica e caracterização estelar, além de ser essencial para o estudo de diversos fenômenos astronômicos. Diante dessas complexidades, emerge a necessidade de uma maior acurácia para estas estimativas.
    
    A determinação da massa de estrelas em sistemas múltiplos pode ser realizada diretamente, mas não representa por completo os aglomerados estelares. Estimar a massa de estrelas em sistemas não-binários envolve frequentemente o uso da relação massa-luminosidade, cujo expoente varia com o estágio evolutivo da estrela. Métodos de ajuste de isócronas também se destacam. Concomitando esses métodos é possível complementar os resultados, dado que cada um tem vantagens e limitações distintas.
    
    A determinação precisa de massas estelares, especialmente em aglomerados jovens e em sistemas não-binários, apresenta desafios significativos devido às limitações inerentes aos métodos convencionais. Em particular, a aplicação direta de relações massa-luminosidade é limitada pela variabilidade dos parâmetros que definem essa relação, dependentes do estágio evolutivo e das propriedades individuais das estrelas \cite{moya2018empirical}.
    
    Esses métodos, baseados em modelos teóricos de evolução estelar, ainda que eficazes para certos tipos estelares, carecem de acurácia para abarcar a diversidade de características presentes em aglomerados de diferentes idades, como aferido por trabalhos como o de \citeonline{serenelli2021weighing}. Assim, a complexidade do problema exige o desenvolvimento de técnicas mais robustas e adaptativas, capazes de integrar múltiplos tipos de dados e metodologias para estimar massas com maior precisão e confiança em uma ampla faixa de condições astrofísicas.
    
    Este estudo tem como objetivo principal a determinação das massas individuais de estrelas de baixa massa em aglomerados jovens, utilizando uma abordagem que integra dados de múltiplas fontes.
    Utilizaremos dois métodos: a relação massa-luminosidade, derivada de modelos teóricos de evolução estelar e o método de ajuste de isócronas. A escolha dos melhores parâmetros da relação massa-luminosidade será feita utilizando técnicas de aprendizado de máquina. Optamos por utilizar dois métodos distintos para que seja possível realizar uma comparação entre os resultados de ambas as metodologias e assim, verificar a consistência e a robustez dos resultados obtidos.
  
    Determinaremos as massas das estrelas dos aglomerados da Nebulosa de Órion (ONC), das Plêiades, de Upper Scorpius (USco), de NGC 869 (h Persei) e de NGC 2516. A escolha desses aglomerados, que representam diferentes estágios evolutivos, visa avaliar a acurácia e as limitações dos métodos aplicados. Nossa escolha também foi pautada pela disponibilidade de períodos de rotação para as estrelas destes aglomerados, uma vez que este trabalho é parte do projeto de pesquisa sobre Rotação Estelar conduzido no Grupo de Formação Estelar da UESC.
